La primera respuesta generada por la IA produjo, precedido de un texto explicativo, un bloque \texttt{flowchart LR} sintácticamente válido, que renderizó sin errores en mermaid.live y respetaba las reglas de la plantilla: una única asociación actor--caso de uso, estereotipos con comillas latinas y \texttt{classDef}/\texttt{class} para el actor. No obstante, tras auditar el código frente a la semántica UML, a los criterios de aceptación de US-8 y al glosario común, se aplicaron las siguientes correcciones y mejoras:
\begin{itemize}[leftmargin=*,nosep]
\item \textbf{Relación «extend» semánticamente incorrecta entre registro e inicio de sesión:} La IA modeló «Iniciar sesión» como extensión de «Registrarse». Iniciar sesión no es un flujo condicional del registro, sino un objetivo propio del cliente. Como la plantilla permite una única asociación actor--caso de uso, se tomó como caso base «Registrarse e iniciar sesión» (título de la historia) y se colgaron de él las funciones opcionales.
\item \textbf{Paso opcional modelado como «include»:} «Guardar direcciones» aparecía como relación «include» de «Registrarse», pero la cliente puede crear su cuenta sin guardar ninguna dirección. Se convirtió en «extend», dirigida desde el caso opcional hacia el base. Se añadió además el caso incluido «Crear cuenta de cliente», paso obligatorio del registro que la IA no representaba.
\item \textbf{Criterios de aceptación y glosario sin respetar:} El criterio habla de un panel con el listado de \emph{pedidos} anteriores, pero la IA usó «Consultar historial de compras», término ausente del glosario común (\texttt{PEDIDO}). Se renombró a «Consultar historial de pedidos» y se dejó como extensión del caso base, igual que «Recordar sesión».
\item \textbf{Actor, identificadores y límite inconsistentes con el proyecto:} La IA usó el actor «Clienta Frecuente» en lugar de \texttt{Cliente} (glosario común), nombres con mayúsculas iniciales en cada palabra, identificadores abreviados como \texttt{ValidarPass} y el límite «E-commerce Boutique de Moda - Módulo de Usuarios». Se unificó el estilo y el límite pasó a titularse «Atelier Jaén - Cuenta de cliente».
\end{itemize}
